% Generated draft: requires mathematical and source review.
\documentclass[11pt]{article}
\usepackage[margin=25mm]{geometry}
\usepackage{fontspec}
\setmainfont{Latin Modern Roman}
\usepackage{amsmath,amssymb,mathtools,unicode-math}
\setmathfont{Latin Modern Math}
\usepackage{enumitem,needspace,xurl,hyperref}
\setlength{\parindent}{0pt}
\setlength{\parskip}{5pt}
\setlength{\emergencystretch}{3em}
\newcommand{\N}{\mathbb{N}}
\newcommand{\Z}{\mathbb{Z}}
\newcommand{\Q}{\mathbb{Q}}
\newcommand{\R}{\mathbb{R}}
\newcommand{\C}{\mathbb{C}}
\newcommand{\F}{\mathbb{F}}
\newcommand{\A}{\mathbb{A}}
\newcommand{\PP}{\mathbb{P}}
\newcommand{\E}{\mathbb{E}}
\newcommand{\eps}{\varepsilon}
\newcommand{\dd}{\,\mathrm{d}}
\begin{document}
\section*{UnsolvedMath ID: 20002825}
\textbf{Source problem code:} catalogue entry
\subsection{Definitions and mathematical statement}
The corrected source asks to describe the graded Hopf isomorphism $SC^{(2)*}\to\Pi^*\cong\mathrm{PiQSym}$ in two presentations and make it compatible with an internal coproduct. The original OCR damages dual stars, indexing, and zero-product conditions. The noncommuting-word presentation uses an alphabet $A=\{a_1,a_2,\ldots\}$ and
\[
 M_\mu=\sum_{\ker(w)=\mu}w,
\]
where $\ker(w)$ partitions positions according to equal letters. This realizes $\Pi=\mathrm{NCSym}$. The matrix-variable/cycle-permutation presentation realizes the relevant dual polynomial object; the extract describes monomials of shape $\prod_i x_{i,\sigma(i)}$ with cycle-support partition $\lambda(\sigma)=\mu$, but does not preserve all exact indexing and quotient relations. The two sides must not silently be treated as the same noncommutative algebra.

Use set partitions $\mathcal P_n$, ordered by refinement ($\nu\le\mu$ means finer), and denote their polynomial basis by $U_\mu$. Let $SC^{(2)}$ be the graded superclass-function Hopf algebra for unitriangular groups over $\F_2$, with superclass indicators $\kappa_\mu$, and take degreewise graded duals. The aid attributes to Aguiar et al. the outer Hopf isomorphism
\[
 \Phi(\kappa_\mu^*)=V_\mu,\qquad
 V_\mu=\sum_{\nu\le\mu}U_\nu.
\]
The full polynomial presentation and external Hopf operations require the cited theorem, not the corrupted display alone.

The aid makes a specific internal-operation choice: transpose pointwise multiplication of superclass functions. Since $\kappa_\mu\kappa_\nu=\delta_{\mu,\nu}\kappa_\mu$, compatibility uniquely requires
\[
 \delta(V_\mu)=V_\mu\otimes V_\mu.
\]
This internal coproduct maps degree $n$ into degree $(n,n)$, unlike the outer Hopf coproduct splitting total degree. With $\mu_{\mathcal P}$ the partition-lattice M\"obius function, its reported transported formula is
\[
 \delta(U_\mu)=\sum_{\nu\le\mu}\mu_{\mathcal P}(\nu,\mu)
 V_\nu\otimes V_\nu.
\]
Equivalently, expanding both factors yields
\[
 \delta(U_\mu)=\sum_{\alpha\vee\beta=\mu}U_\alpha\otimes U_\beta,
\]
where join is in $\mathcal P_n$. Indeed the coefficient is the finite M\"obius sum over common upper bounds below $\mu$. This formula specifies the $U$ expansion without assuming those elements are group-like.

For $\mu=12$, the one-block partition of two elements, the source's explicit test is
\[
 \delta(U_{12})=U_{1|2}\otimes U_{12}
 +U_{12}\otimes U_{1|2}+U_{12}\otimes U_{12}.
\]
Thus $U_{12}$ is not group-like, while $V_{12}$ is. The conclusion is conditional on the pointwise-dual interpretation of ``internal''; choosing a different internal operation is a different compatibility question. A complete identification with both source presentations still requires the missing exact polynomial conventions and authoritative reference.
\subsection{Short English statement}
The aid's outer isomorphism sends a dual superclass indicator to a sum of polynomial basis elements. Transposing pointwise multiplication makes those sums group-like and forces a Möbius/join formula for the individual basis elements. The degree-two cross terms rule out the naive group-like rule for the $U$ basis.
\subsection{Sources}
\begin{itemize}
\item Local UnsolvedMath registry, permanent ID 20002825, source code catalogue entry; created and updated 14 August 2026. Canonical \texttt{aim-representation-theory-notes.json}, zero-based \texttt{notes[96]}, tag \texttt{problem}, item 1.13.
\item AIM ``Supercharacters and combinatorial Hopf algebras'': \url{https://aimath.org/WWN/supercharacters/supercharacters.pdf}. Statement status is corrected\_verified according to the supplied verification, which reports PDF comparison with the later polynomial realization. The original OCR includes a page footer and damaged quotient conditions; original remarks and literature field are empty.
\item Research attempt 1: \texttt{partial\_result}, problem status at run \texttt{partially\_solved}, rigor label \texttt{proved}, novelty confidence low. It cites Aguiar et al. [AAB] without a full author list, title, publication data, or complete theorem statement.
\item Ulam AI / UnsolvedMath Contributors, CC BY 4.0. Machine-generated research aid, not peer reviewed; independently verify before citation.
\end{itemize}
\end{document}
